\setvariables[meta]
  [title={Huffman Kodierung},
   subtitle={Zusammenfassung der Funktionsweise},
   date={2024-10-21},
   author={Niklas von Hirschfeld, Finn Garbers},
  ]

\setupinteraction[title={\getvariable{meta}{title}}]

% \showlayout
% \showframe

\setuplayout[
   backspace=1cm, % Space on the left margin
   width=fit, % Adjust width to fit remaining space
   height=fit,
   topspace=0.5cm, % Top margin
   header=0.5cm, % Header size
   footer=1cm, % Footer size
   bottomspace=1cm, % Bottom margin
   leftedge=1.5cm,
   rightedge=1cm,
   rightmargin=3.5cm, % Space for margin notes
   rightmargindistance=0.5cm
]

\setuppagenumbering[location=]
\setupheader[text][before={},after={}]
\setupheadertexts[][]

\setuphead[chapter]
          [
			header=empty,
			align=middle,
            style=\itd,
            number=no,
            before={\blank[big,force]},
            after={\blank[big,force]},
          ]

\starttext
\startbodymatter

\chapter{Huffman Codierung}

\startframedtext[location=middle,align=middle,width=\dimexpr\textwidth\relax,frame=off]
Ein Algorithmus zur {\bf verlustfreien} Datenkompression, der von {\bf David A. Huffman} in
{\bf 1952} entwickelt wurde. Sie erstellt kürzere Codes für häufiger auftretende
Zeichen, während seltenere Zeichen längere Codes erhalten. Ziel ist es, die
Größe der zu speichernden oder zu übertragenden Daten zu reduzieren um somit
die Übertragungsdauer zu verringern.
\stopframedtext


\section{Funktionsweise}

Jedem einzigartigen {\bf Quellsymbol} wird ein festes {\bf Codewort}
zuwgewiesen. Die Codewörte variieren in der länge. Um möglichst effizient die
Codewörter zu gestalten, wird dem Quellsymbol, welches am häufigsten genutzt
wird auch das kürzte Codewort zugewiesen.

Dabei ist es wichtig, das die Symbole {\bf Präfixfrei}\sidenote{{\em
Präfixfrei} heißt, kein Symbol darf mit einem anderen beginnen. So ist folgende
Symbolgruppe Präfixfrei: $\{0, 100, 101, 11\}$, aber $\{10, 100, 101, 11\}$
nicht.} sind. Dies ist nötig, um Symbole unterschiedlicher länge eindeutig zu erkennen.

\section{Aufbau eines Baumes}

Um korrekte Codewörter zu finden, ist die gängiste möglichkeit über einen
$k$-nären Baum\sidenote{ein Baum desse Knoten jeweils $k$-viele Blätter
aufweisen}. Folgend wird die vorgehensweise mittels eines binären Baums
erläutert:

\startitemize[packed,m]
\item Für jedes Zeichen wird die Häufigkeit bestimmt
\item Aufbauen des Binärbaum\sidenote{Ein {\bf Binärbaum} ist ein Baum, bei dem jeder Knotem {\bf maximal zwei}}
\startitemize[packed,m]
\item Die zwei kleinsten Knoten\sidenote{Zu beginn wird jeder Buchstabe als eigener Knoten betrachtet} finden (Falls es mehr als zwei Knoten gibt, wird der Knoten mit der geringsten Tiefe gewählt, ansonsten kann frei gefählt werden)
\item Diese mit einem Knoten verknüpfen und die Häufigkeit der beiden Blätter addiert in den Knoten schreiben, bzw. speichern.
\item Diese Schritte wiederholen, bis ein einzelner Binärbaum übrig ist.
\stopitemize
\item Nun wird an jeden linken Ast eine Eins geschrieben und an jeden rechten eine 0\sidenote{Es können natürlich beliebiege Zeichen gewählt werden, $0$ und $1$ bieten sich bei der speicherung auf einem Computer an, da sie ohne weiter Codierung direkt gespeichert werden können.}
\stopitemize

\section{Beispiel: Mississippi}

\startbuffer[beispiel]
\startMPpage

u := 1.5cm;

pair M,I,S,P, MP, SMP, SMPI;

SMPI=(0, 0);
SMP=(-1 * u, -1 * u);
I=(1 * u, -1 * u);
S=(-2 * u, -2 * u);
MP=(0, -2 * u);
P=(1 * u, -3 * u);
M=(-1 * u, -3 * u);

draw M--MP;
draw P--MP;

draw S--SMP;
draw MP--SMP;

draw SMP--SMPI;
draw I--SMPI;

color lightyellow;
lightyellow := (1, 1, 0.8);

fill fullcircle scaled u shifted I withcolor lightyellow;
draw fullcircle scaled u shifted I;
fill fullcircle scaled u shifted S withcolor lightyellow;
draw fullcircle scaled u shifted S;
fill fullcircle scaled u shifted P withcolor lightyellow;
draw fullcircle scaled u shifted P;
fill fullcircle scaled u shifted M withcolor lightyellow;
draw fullcircle scaled u shifted M;
fill fullcircle scaled u shifted MP withcolor white;
draw fullcircle scaled u shifted MP;
fill fullcircle scaled u shifted SMP withcolor white;
draw fullcircle scaled u shifted SMP;
fill fullcircle scaled u shifted SMPI withcolor white;
draw fullcircle scaled u shifted SMPI;

label("I", I);
label("S", S);
label("P", P);
label("M", M);
label("S", S);
label("MP", MP);
label("SMP", SMP);
label("SMPI", SMPI);

% label.ulft("1", 0.5[SMP--SMPI]);

label.rt("M: 1", (-4 * u, 0));
label.rt("I: 4", (-4 * u, -0.5 * u));
label.rt("S: 4", (-4 * u, -1 * u));
label.rt("P: 2", (-4 * u, -1.5 * u));

\stopMPpage
\stopbuffer

\placefigure[force][beispiel]{Beispiel der Huffmancodierung mittels eines Binärbaumes.}{
\typesetbuffer[beispiel]
}

\stopbodymatter

\stoptext
